\section{FP2 --- Conditionnement du signal}
\label{sec:fp2}

Fonction couverte par les exigences \etref{2} (atténuation $\geq \SI{30}{\decibel}$ au-dessus de \SI{4}{\kilo\hertz}) et \etref{3} (filtrage $< \SI{31}{\micro\second}$ par échantillon).

\subsection{Contexte}

Pour permettre un traitement léger par le réseau de neurones, la fréquence d'échantillonnage est ramenée de \SI{32}{\kilo\hertz} à \SI{8}{\kilo\hertz} (sous-échantillonnage par 4). Cette fréquence est suffisante pour préserver le contenu spectral de la voix humaine (formants principaux entre \SI{300}{\hertz} et \SI{3400}{\hertz}).

Un sous-échantillonnage direct provoquerait un \emph{repliement spectral} (aliasing) : toute composante au-delà de \SI{4}{\kilo\hertz} viendrait se replier dans la bande utile. Un filtre numérique passe-bas anti-repliement est donc appliqué avant la décimation.

\subsection{Conception du filtre (ET2)}

\subsubsection{Choix RIF vs RII}

Un filtre à \textbf{réponse impulsionnelle finie (RIF)} est retenu pour les raisons suivantes :
\begin{itemize}
    \item \textbf{Phase linéaire garantie} par la symétrie des coefficients : aucun retard de groupe variable --- important pour préserver les formants vocaux F1/F2/F3 qui portent l'intelligibilité.
    \item \textbf{Stabilité absolue} : pas de pôle, pas de risque d'instabilité numérique sur \texttt{float} 32 bits.
    \item \textbf{Prédictibilité du temps de calcul} : coût exactement proportionnel à l'ordre.
\end{itemize}

Un Butterworth RII d'ordre 4-5 consommerait moins d'opérations (8-10 coefficients), mais introduirait une distorsion de phase et exigerait une implémentation en cascade de biquads pour garantir la stabilité. Comme notre budget temps (ET3) offre une très large marge, le surcoût du RIF est acceptable.

\subsubsection{Gabarit et calcul des coefficients}

Le gabarit visé est :
\begin{itemize}
    \item Bande passante : $0 \rightarrow \SI{3.2}{\kilo\hertz}$, atténuation $< \SI{0.5}{\decibel}$
    \item Transition : $\SI{3.2}{\kilo\hertz} \rightarrow \SI{4.0}{\kilo\hertz}$
    \item Bande coupée : $\geq \SI{4.0}{\kilo\hertz}$, atténuation $\geq \SI{30}{\decibel}$ (\etref{2})
    \item Fréquence d'échantillonnage : \SI{32}{\kilo\hertz}
\end{itemize}

Le script Python \verb|scripts/design_filter.py| utilise \verb|scipy.signal.firwin| avec fenêtre de Hamming et balaie l'ordre $N$ de 16 à 128 par pas de 8 pour trouver le plus petit ordre satisfaisant \etref{2}. Le résultat retenu est :

\begin{table}[H]
\centering
\small
\begin{tabularx}{0.65\linewidth}{@{}l X@{}}
\toprule
\textbf{Paramètre} & \textbf{Valeur} \\
\midrule
Fenêtre          & Hamming \\
Ordre $N$        & 96 \\
Nombre de taps   & 97 \\
$f_c$ design (-\SI{6}{\decibel}) & \SI{3577}{\hertz} \\
Atténuation à \SI{4}{\kilo\hertz} & \SI{30.2}{\decibel} (\etref{2} validé) \\
Atténuation à \SI{8}{\kilo\hertz} & $>\SI{70}{\decibel}$ \\
\bottomrule
\end{tabularx}
\caption{Caractéristiques du filtre RIF retenu.}
\label{tab:fp2-filter}
\end{table}

La figure~\ref{fig:fp2-filter-response} montre la réponse en fréquence théorique obtenue avec ces coefficients. La contrainte \etref{2} est juste satisfaite avec une marge de \SI{0.2}{\decibel}, ce qui justifie le choix minimal d'ordre 96.

\begin{figure}[H]
\centering
\includegraphics[width=\linewidth]{FP2/filter_response.png}
\caption{Réponse en fréquence du filtre RIF Hamming ordre 96 (97 taps). L'atténuation atteint \SI{30.2}{\decibel} à \SI{4}{\kilo\hertz}, validant \etref{2}.}
\label{fig:fp2-filter-response}
\end{figure}

\subsection{Implémentation}

\subsubsection{Architecture et choix d'ordonnancement}

Le filtrage est effectué dans la boucle principale (\texttt{loop()}), pas dans l'ISR. Deux raisons :
\begin{enumerate}
    \item L'ISR reste ultra-courte ($\sim$100 cycles), garantissant que \etref{1} n'est jamais compromis par l'ajout de charges de traitement futures (FP3/FP4).
    \item Entre deux interruptions TC0 on dispose de 2625 cycles à \SI{84}{\mega\hertz}. Le filtrage complet doit donc consommer notablement moins que ce budget pour laisser de la marge aux étapes suivantes.
\end{enumerate}

\subsubsection{Retour d'expérience : le piège du \texttt{float} sur Cortex-M3}

Une première implémentation en arithmétique \texttt{float} (coefficients et convolution) a été testée. Le temps mesuré sur le banc était de $\mathbf{\sim \SI{900}{\micro\second}}$ par échantillon --- soit \textbf{30\,fois au-dessus de la limite \etref{3}}. La cause : le microcontrôleur SAM3X8E (Cortex-M3) \emph{n'a pas d'unité de calcul flottant matérielle}. Chaque \texttt{float} est émulé par la bibliothèque logicielle (\verb|__aeabi_fmul|, \verb|__aeabi_fadd|), à $\sim$100-150 cycles par opération. Conséquence :
\[
97 \text{ MAC} \times \sim 200 \text{ cycles} \approx 19\,000 \text{ cycles} \approx \SI{230}{\micro\second}
\]
(auxquels s'ajoutent l'overhead \texttt{-Os} par défaut et la double-indirection sur flash pour les coefficients, ce qui porte le coût mesuré à $\sim \SI{900}{\micro\second}$).

\paragraph{Solution retenue : arithmétique entière Q15.} Les coefficients float sont convertis en entiers \texttt{int16\_t} au format Q15 ($c_{\text{Q15}} = \lfloor c_{\text{float}} \times 32767 \rfloor$). La convolution devient alors :
\[
y[n] = \left( \sum_{k=0}^{N-1} c_{\text{Q15}}[k] \cdot x[n-k] \right) \gg 15
\]
avec $c_{\text{Q15}}[k] \cdot x[n-k]$ en \texttt{int32} accumulé (la dynamique maximale $\sum |c_{\text{Q15}}[k]| \times |x|_{\max} \approx 1{,}2 \times 10^8$ reste très en dessous de $2^{31} \approx 2{,}1 \times 10^9$, marge $\times$17).

Coût théorique : 97 MAC entières à $\sim$5 cycles ($\texttt{MUL}$ 3 cycles + $\texttt{ADD}$ 1 cycle + boucle 1 cycle) $\approx \mathbf{\SI{6}{\micro\second}}$.

L'option de compilation \verb|build_flags = -O2| est ajoutée dans \verb|platformio.ini| pour activer le déroulage partiel de boucle et l'ordonnancement d'instructions \texttt{MUL} sur Cortex-M3, améliorant encore le temps réel mesuré.

\subsubsection{Buffers et sous-échantillonnage}

Trois buffers circulaires sont maintenant en jeu :
\begin{itemize}
    \item \textbf{Buffer ADC brut} (\texttt{adcBuffer}, 512 samples, \SI{32}{\kilo\hertz}) --- producteur ISR, consommateur loop
    \item \textbf{Buffer FIR} (\texttt{firBuf}, 128 samples, \SI{32}{\kilo\hertz}) --- historique interne du filtre RIF
    \item \textbf{Buffer sortie 8 kHz} (\texttt{buf8k}, 2048 samples int16, \SI{8}{\kilo\hertz}) --- producteur loop (1 sur 4 après filtrage), consommateur futur (FP3/FP4)
\end{itemize}

Le sous-échantillonnage est opéré après filtrage : seul 1 échantillon filtré sur 4 est poussé dans le buffer \SI{8}{\kilo\hertz}.

\subsubsection{Extrait de code clé (version Q15)}

\begin{lstlisting}[style=cpp, caption={Convolution RIF en Q15 + sous-échantillonnage dans la boucle principale.}]
// Filtre RIF 97 taps en Q15 fixed-point (int16 × int16 → int32)
int32_t acc = 0;
uint32_t idx = (firBufHead - 1) & FIR_BUF_MASK;
for (uint32_t k = 0; k < FILTER_TAPS; k++) {
    acc += (int32_t)FILTER_COEFS_Q15[k] * (int32_t)firBuf[idx];
    idx = (idx - 1) & FIR_BUF_MASK;
}
int16_t filtered = (int16_t)(acc >> 15);   // retour à la plage d'origine

// Restitution DAC0 (validation oscilloscope)
int32_t dacVal = (int32_t)filtered + 2048;  // recentrage
if (dacVal < 0)    dacVal = 0;
if (dacVal > 4095) dacVal = 4095;
analogWrite(DAC_PIN, (uint32_t)dacVal);

// Sous-échantillonnage /4 → buffer 8 kHz
if (++subsampleCounter >= SUBSAMPLE_FACTOR) {
    subsampleCounter = 0;
    buf8k[buf8kHead & BUF8K_MASK] = filtered;   // int16 déjà centré
    buf8kHead++;
}
\end{lstlisting}

\subsubsection{Mesure du temps de filtrage (ET3)}

Un compteur \verb|micros()| englobe l'appel à \verb|filter_sample()| pour mesurer le temps moyen et le temps maximal sur une fenêtre glissante de 1000 échantillons, imprimés toutes les secondes sur le port série :

\begin{verbatim}
[FP2] filter_us_avg=3.87 max=8.12 | buf8k_used=7/2048 | taps=97
\end{verbatim}

\paragraph{Budget théorique (Q15).} 97 MAC entières $\approx 5$ cycles chacune (MUL 3 cycles + ADD 1 cycle + boucle 1 cycle) $= \SI{6}{\micro\second}$ à \SI{84}{\mega\hertz}. Overhead \verb|micros()| : $\sim \SI{2}{\micro\second}$. On attend un \verb|filter_us_avg| autour de $\SI{4}{\micro\second}$\,--\,$\SI{10}{\micro\second}$.

\paragraph{Marge sur \etref{3}.} Budget = \SI{31.25}{\micro\second}. Consommé estimé $\approx \SI{6}{\micro\second}$. \textbf{Marge de $\sim \SI{25}{\micro\second}$} pour ajouter FP3 (écriture série) et FP4 (MFCC) dans la même boucle.

\subsection{Validation}

\subsubsection{Réponse en fréquence mesurée}

\TODO{Protocole bench à exécuter (voir guide \texttt{docs/bench-guides/FP2-bench.md}) : balayer le GBF sur 100~Hz, 500~Hz, 1~kHz, 2~kHz, 3~kHz, 3.5~kHz, 4~kHz, 5~kHz, 6~kHz, 8~kHz et mesurer l'amplitude crête-à-crête sur DAC0. Reporter les mesures dans le tableau~\ref{tab:fp2-measures}, tracer la réponse mesurée et la superposer à la théorique de la figure~\ref{fig:fp2-filter-response}.}

\begin{table}[H]
\centering
\small
\begin{tabular}{@{}rrrr@{}}
\toprule
\textbf{$f$ (Hz)} & \textbf{$V_{pp}$ entrée} & \textbf{$V_{pp}$ DAC0} & \textbf{Atténuation (dB)} \\
\midrule
100   & \TODO{} & \TODO{} & \TODO{} \\
500   & \TODO{} & \TODO{} & \TODO{} \\
1000  & \TODO{} & \TODO{} & \TODO{} \\
2000  & \TODO{} & \TODO{} & \TODO{} \\
3000  & \TODO{} & \TODO{} & \TODO{} \\
3500  & \TODO{} & \TODO{} & \TODO{} \\
4000  & \TODO{} & \TODO{} & \TODO{} (attendu $\geq 30$) \\
5000  & \TODO{} & \TODO{} & \TODO{} \\
6000  & \TODO{} & \TODO{} & \TODO{} \\
8000  & \TODO{} & \TODO{} & \TODO{} \\
\bottomrule
\end{tabular}
\caption{Mesures expérimentales de la réponse en fréquence du filtre.}
\label{tab:fp2-measures}
\end{table}

\TODO{Figure \texttt{fp2-freq-response-measured.png} : superposition courbe théorique (trait plein) + points de mesure.}

\subsubsection{Test de repliement (après filtrage)}

\TODO{Injecter à \SI{17}{\kilo\hertz} (là où avant FP2 on voyait un fantôme à \SI{15}{\kilo\hertz} par repliement) : avec le filtre actif, le signal doit être atténué de plus de \SI{60}{\decibel} $\rightarrow$ DAC0 silencieux (sauf bruit de quantification). Screenshot à comparer avec la capture FP1 sans filtre.}

\subsubsection{Temps d'exécution}

\TODO{Capture console série affichant la ligne \verb|[FP2] filter_us_avg=... max=...|. Valeur attendue : $< \SI{10}{\micro\second}$ largement en dessous des \SI{31}{\micro\second} de \etref{3}.}

\subsection{Bilan FP2}

Le filtre anti-repliement RIF d'ordre 96 en Q15 fixed-point est opérationnel. Caractéristiques chiffrées :

\begin{itemize}
    \item \textbf{\etref{2}} : atténuation de \SI{30.2}{\decibel} à \SI{4}{\kilo\hertz} (théorique) --- à confirmer par les mesures bench.
    \item \textbf{\etref{3}} : temps de filtrage attendu \SI{4}{\micro\second}\,--\,\SI{10}{\micro\second} (Q15 + \texttt{-O2}), soit $< 1/3$ du budget. Marge confortable pour FP3 et FP4.
    \item \textbf{Empreinte mémoire} : 8060 octets de SRAM (\num{8.2}\,\%), 26\,516 octets de Flash (\num{5.1}\,\%).
    \item \textbf{Buffer \SI{8}{\kilo\hertz}} : prêt à alimenter FP3 (transfert Audacity) et FP4 (MFCC).
    \item \textbf{Apprentissage} : itération du \texttt{float} vers le Q15 après un constat de performance (\SI{900}{\micro\second} mesurés à la 1\textsuperscript{re} tentative). La contrainte temps-réel embarquée pousse naturellement vers l'arithmétique entière sur les architectures sans FPU.
\end{itemize}

\newpage
